323111047

\section{Introducción}

Planteamiento del problema. El reto técnico de esta actividad consiste en desarrollar una aplicación de consola en lenguaje C que implemente el criptosistema de Vigenère, el cual es un cifrador por sustitución polialfabética. El problema central radica en gestionar un alfabeto de 27 letras incluyendo la ñ y aplicar fórmulas de aritmética modular para realizar tanto el cifrado como el descifrado de mensajes. Para la resolución de este problema, es indispensable el uso de estructuras de datos struct y el manejo de acceso a miembros mediante el operador de desreferencia flecha, asegurando que los caracteres no alfabéticos se mantengan intactos durante el proceso.

\vspace{0.4cm}

Motivación. La importancia de resolver este problema se basa en la necesidad de comprender métodos de protección de información más robustos que los de sustitución simple, como el cifrado César. Al utilizar una clave que varía la posición del desplazamiento para cada letra, el método de Vigenère ofrece una mayor seguridad contra el análisis de frecuencias básico. Estudiar este algoritmo permite al estudiante de ingeniería entender cómo la lógica matemática aplicada a la programación puede fortalecer la privacidad de los datos, sirviendo como fundamento para comprender sistemas criptográficos más avanzados utilizados en la actualidad.

\vspace{0.4cm}

Objetivos. El propósito fundamental de esta práctica es lograr la implementación funcional de un menú que permita al usuario cifrar y descifrar textos de manera interactiva. Específicamente, se busca aplicar correctamente el operador flecha para la manipulación de datos dentro de estructuras y validar que el programa respete el alfabeto castellano de 27 caracteres en todas sus operaciones. Finalmente, se pretende verificar que el uso de la aritmética modular garantice la integridad del mensaje al realizar la operación inversa de descifrado, demostrando así la eficiencia del algoritmo polialfabético implementado.